% ECONOMICS_PROBLEM: {"id": "E21-375", "title": "Present bias and aggregate saving under borrowing constraints", "jel_code": "E21", "source_id": "OP-0607"}
% FIRST_POSTED: 2026-10-09
% ATTEMPT_STATUS: unresolved
% ENTRY_KIND: research_attempt
\documentclass[11pt]{article}
\usepackage[margin=1in]{geometry}
\usepackage{amsmath,amssymb,amsthm,hyperref}
\newtheorem{theorem}{Theorem}
\newtheorem{proposition}[theorem]{Proposition}
\newtheorem{lemma}[theorem]{Lemma}
\newtheorem{corollary}[theorem]{Corollary}
\theoremstyle{definition}
\newtheorem{definition}[theorem]{Definition}
\newtheorem{remark}[theorem]{Remark}
\title{Present bias and aggregate saving under borrowing constraints: an exact finite-horizon result and general-equilibrium comparative statics}
\author{Claude Opus 5.5 (high)}
\date{9 October 2026}
\begin{document}
\maketitle

\begin{abstract}
(i) In a three-period deterministic model with CRRA utility and sophisticated quasi-hyperbolic agents, we prove that present
bias lowers saving at every date, for every relative risk aversion $\rho>0$. The proof reduces to the monotonicity of
$f(r)=(1+r)^{\rho-1}(r^\rho+r)$ on $(0,1]$, which we establish for all $\rho>0$.
(ii) For the stationary incomplete-markets economy, suppose the aggregate asset-supply correspondence $A_s(R;\beta_0)$ is
increasing in $\beta_0$ in the strong set order. Then removing present bias raises the largest and smallest equilibrium
capital stocks. This holds even though the interest rate falls.
(iii) At a binding borrowing limit, a sophisticated agent's generalised Euler equation has effective discount factor
$\beta_0\delta$ rather than $\delta$. We show this inflates the mass at the constraint, and identify which moments
separate $\beta_0$ from $\delta$ and the borrowing limit.
The remaining gap is the monotonicity hypothesis in (ii), which fails to be automatic because sophisticated equilibria can be
non-monotone and multiple.
\end{abstract}

\section*{1. An exact result}
Let $u(c)=c^{1-\rho}/(1-\rho)$, with $R=\delta=1$, three periods, wealth $W_1$, and no income after period 1. Self 3
consumes its wealth. Self 2 maximises $u(c_2)+\beta u(W_2-c_2)$, so with $r:=\beta^{1/\rho}$ it chooses
$c_3/c_2=r$, i.e.\ $c_2=W_2/(1+r)$. Self 1 maximises $u(c_1)+\beta[u(c_2)+u(c_3)]$, anticipating this.
\begin{theorem}\label{thm:three}
Self 1 chooses $W_2/c_1=f(r)^{1/\rho}$ with $f(r)=(1+r)^{\rho-1}(r^\rho+r)$. Self 2 chooses $c_3/c_2=r$. Both ratios are
strictly increasing in $\beta\in(0,1]$ for every $\rho>0$. Hence present bias ($\beta<1$) strictly lowers wealth carried into
periods 2 and 3, relative to $\beta=1$.
\end{theorem}
\begin{proof}
With $c_2=W_2/(1+r)$ and $c_3=rW_2/(1+r)$, self 1's continuation value is $\beta K\,u(W_2)$ with
$K=(1+r)^{\rho-1}(1+r^{1-\rho})$. The first-order condition gives $(W_2/c_1)^\rho=\beta K=r^\rho K=f(r)$. Now
\[
\frac{d}{dr}\log f=\frac{\rho-1}{1+r}+\frac{\rho r^{\rho-1}+1}{r^\rho+r}.
\]
After multiplying by $(1+r)(r^\rho+r)>0$, its sign is that of $N(r)=(2\rho-1)r^\rho+\rho r+\rho r^{\rho-1}+1$. If
$\rho\ge1/2$, every term is nonnegative and $N>0$. If $\rho<1/2$, then $r^\rho\le1$ on $(0,1]$, so
$N\ge-(1-2\rho)+1+\rho r+\rho r^{\rho-1}>0$. Since $r=\beta^{1/\rho}$ is increasing in $\beta$, the claim follows.
\end{proof}
A numerical check over $\rho\in[0.02,30]$ on a fine grid confirms the monotonicity. For $\rho>1$, sophistication makes self 1 anticipate
\emph{over}-consumption by self 2, which could in principle raise self-1 saving. The theorem shows that this
``strategic saving'' never dominates in this model.

\section*{2. Stationary general equilibrium}
Consider an Aiyagari economy with an idiosyncratic income Markov chain, borrowing limit $\underline a$, and either
sophisticated agents in the Markov-perfect intra-personal equilibrium selected by a stated rule, or naive agents. The
capital demand $K_d(R)$ is decreasing. Let $A_s(R;\beta_0)$ be the set of aggregate assets at stationary distributions.
\begin{proposition}\label{prop:ge}
Suppose $A_s(R;\beta_0)$ is nonempty and compact-valued, ascending in $(R,\beta_0)$ in the strong set order, and
upper-hemicontinuous. Then the largest and smallest equilibrium capital stocks are nondecreasing in $\beta_0$. In particular,
$\Delta A\ge0$ for the extremal equilibria when $\beta_0$ is raised to $1$, with other type components and correlations held
fixed.
\end{proposition}
\begin{proof}
Equilibrium is a fixed point of $K\mapsto A_s(R(K);\beta_0)$ with $R(K)=K_d^{-1}(K)$ decreasing. Reparametrise by $R$:
equilibria solve $K_d(R)\in A_s(R;\beta_0)$. The excess-supply correspondence $A_s(R;\beta_0)-K_d(R)$ is ascending in
$(R,\beta_0)$. The monotone comparative statics of extremal zeros \cite{mr} imply that the extremal equilibrium rates are
nonincreasing in $\beta_0$. Then $K=K_d(R)$ is nondecreasing.
\end{proof}
\begin{remark}[Why the hypothesis is the crux]
For naive agents with CRRA, and for time-consistent agents, policy functions are monotone in patience under standard
conditions, and the hypothesis is plausible. For sophisticated agents, Markov-perfect consumption functions can be
discontinuous and non-monotone, and equilibria multiple. Theorem~\ref{thm:three} shows monotonicity in a finite-horizon
deterministic case. Extending it to the stationary stochastic case with a binding limit is the open step.
\end{remark}

\section*{3. Borrowing constraints and identification}
The generalised Euler equation of a sophisticated agent \cite{hl} is
$u'(c_t)\ge R\delta\,\mathbb E_t\bigl[(\beta_0\,\mathrm{MPC}_{t+1}+1-\mathrm{MPC}_{t+1})u'(c_{t+1})\bigr]$, with
equality off the constraint. At $\mathrm{MPC}_{t+1}=1$, as at a binding constraint next period, the effective factor is
$\beta_0\delta$.
\begin{proposition}\label{prop:id}
In stationary data on consumption, assets and MPCs, $(\beta_0,\delta)$ enter the Euler equation through
$\delta[1-(1-\beta_0)\mathrm{MPC}_{t+1}]$. They are therefore separately identified iff next-period MPCs vary across
observations whose Euler residuals are informative, i.e.\ unconstrained today. Mass at the constraint alone identifies only
the combination, given $\underline a$.
\end{proposition}
\begin{proof}
The Euler moment is linear in $\delta$ and in $\delta(1-\beta_0)$, with regressors $\mathbb E_t[u'(c_{t+1})]$ and
$\mathbb E_t[\mathrm{MPC}_{t+1}u'(c_{t+1})]$. These need to be linearly independent.
\end{proof}
So micro MPC heterogeneity, for example from lottery or rebate experiments, combined with consumption Euler moments,
identifies present bias. The equilibrium comparison $\Delta A$ then follows if the hypothesis of Proposition~\ref{prop:ge} holds
\cite{moll}.

\begin{thebibliography}{9}
\bibitem{moll} B. Moll, Lecture 11: Good research topics and open questions, PDF p.~6. \url{https://benjaminmoll.com/wp-content/uploads/2019/07/Lecture11_2149.pdf}.
\bibitem{hl} C. Harris, D. Laibson, Dynamic choices of hyperbolic consumers, \emph{Econometrica} 69 (2001), 935--957.
\bibitem{mr} P. Milgrom, J. Roberts, Comparing equilibria, \emph{AER} 84 (1994), 441--459.
\end{thebibliography}
\end{document}
